% Paper 29: The GeoVac Hopf Graphs Are Ramanujan
% GeoVac Project, April 2026
% v1.1 (Sprint-2 follow-up sync: RH-D bound-crossing, RH-E Hypothesis-1 validation, RH-G spectral-zeta redirection)

\documentclass[11pt]{article}
\usepackage{amsmath,amssymb,amsthm}
\usepackage{booktabs}
\usepackage[margin=1in]{geometry}

\newtheorem{theorem}{Theorem}
\newtheorem{proposition}{Proposition}
\newtheorem{corollary}{Corollary}
\newtheorem{definition}{Definition}
\newtheorem{observation}{Observation}
\newtheorem{hypothesis}{Hypothesis}

\title{The GeoVac Hopf Graphs Are Ramanujan:\\
Ihara Zeta, Graph-RH, and a Scope Boundary on Selberg-on-Hydrogen}
\author{GeoVac Project}
\date{April 2026}

\begin{document}
\maketitle

\begin{abstract}
We compute the Ihara zeta function of the two GeoVac Hopf graphs -- the
Fock-projected $S^3$ graph for the Coulomb problem~\cite{paper7} and the
Bargmann--Segal $S^5$ graph for the 3D isotropic harmonic
oscillator~\cite{paper24} -- via the Ihara--Bass determinantal
formula~\cite{ihara1966,bass1992} and the Hashimoto non-backtracking
edge operator~\cite{hashimoto1989}.  Four results emerge.
(1)~\emph{Graph-Ramanujan at small sizes, bound-crossing at modest
sizes}: every graph tested at small sizes
($S^3$ at $n_{\max}=2,3$; $S^5$ at $N_{\max}=2,3$) satisfies the
Kotani--Sunada~\cite{kotani2000} graph-RH (Ramanujan) bound
\emph{strictly}, with deviations of the largest non-trivial Hashimoto
eigenvalue from $\sqrt{q_{\max}}$ ranging from $-0.198$ ($S^3$,
$n_{\max}=3$) to $-0.357$ ($S^5$, $N_{\max}=3$).  A follow-on
post-draft extension to $n_{\max},N_{\max}\in\{4,5\}$ establishes
that three of four graph families \emph{cross} the Kotani--Sunada
bound at modest sizes ($V\approx 30$--$60$), with the signed
deviation growing linearly in $V$; the graph-RH for GeoVac is
therefore a \emph{finite-size} statement, not an asymptotic one.
(2)~\emph{Algebraicity of zeros}:
every $\zeta_G(s)^{-1}$ factors as a product of integer-coefficient
polynomials; no transcendental ($\pi$, $\zeta(2)$, $\zeta(3)$) enters
the zero set.  This sharpens Paper 24's $\pi$-free
certificate~\cite{paper24} from the adjacency data to the Ihara zeta.
(3)~\emph{Structural factorization, validated}: the $S^3$ Coulomb zeta
decomposes per $\ell$-shell because the Fock adjacency preserves~$\ell$;
the $S^5$ $N_{\max}=3$ zeta factors as $(s\pm 1)^{23}\cdot
P_{12}(s^2)\cdot P_{22}(s^2)$, and the $12+22$ dichotomy is \emph{exactly}
the $\mathbb{Z}_2$ $m_\ell\to-m_\ell$ block decomposition of the
node-level Ihara--Bass matrix --- the only sub-action of the
Paper~25~\cite{paper25} Hopf $U(1)$ that commutes with a real integer
adjacency.  The analogous $22+24$ decomposition holds for the
Dirac-$S^3$ Rule B $n_{\max}=3$ graph under $m_j\to-m_j$.
(4)~\emph{Scope boundary}: Fock's Wick rotation $S^3\to H^3$ is
mathematically clean~\cite{bander1966} but produces no natural discrete
quotient $H^3/\Gamma$ from framework invariants; the Selberg-zeta
direction on hyperbolic quotients of hydrogen is therefore structurally
closed.  The Ihara-to-classical-$\zeta$ bridge is additionally blocked
by bipartiteness (odd-argument Dirichlet characters such as $\chi_{-4}$
have support disjoint from the Ihara walk-length spectrum); the
spectral-side Dirichlet series of Paper 28 remains the open direction.
\end{abstract}

% ===================================================================
\section{Introduction}
% ===================================================================

The GeoVac framework~\cite{paper1,paper7} computes quantum-mechanical
observables from discrete graph Laplacians on compact manifolds.  In
the accumulated results of the project, six transcendental constants
have appeared at positive integer arguments of the Riemann zeta and
Dirichlet $L$ family:
$\pi$ (calibration~\cite{paper18});
$\zeta(2)=\pi^2/6$ as a Dirichlet series of the Fock degeneracy lattice
at the packing exponent $d_{\max}=4$;
$\zeta(3)$ in the Dirac sector at $s=d_{\max}$ and in flat-space QED
sums~\cite{paper28};
Catalan's constant $G=\beta(2)$ and the Dirichlet $L$-value $\beta(4)$
from two-loop vertex parity on $S^3$~\cite{paper28}; and
the irreducible $S_{\min}\approx 2.47954$~\cite{paper28}.
Every member of this tally lives at a positive integer argument of
$\zeta$: they are \emph{periods}, not \emph{zeros}.  The Riemann
Hypothesis is about the location of nontrivial zeros of $\zeta$ on
the critical line $\Re s = \tfrac12$.  Periods and zeros are
structurally different objects; the standard bridges (explicit
formula, Euler product) move data across the divide but do not
collapse the distinction.

This paper reports what happens when one approaches the RH-bridge
question from inside GeoVac rather than from outside.  We do not
attempt to prove or disprove RH.  We ask: does the framework contain,
on its own, a graph-zeta or spectral-zeta object with an RH-type
structure --- an Euler product, a functional equation, and a critical
locus that our zeros respect?  The answer, for the two Hopf graphs
the framework natively builds, is yes.  Every tested GeoVac graph
satisfies the Kotani--Sunada~\cite{kotani2000} Ramanujan bound
(Section~\ref{sec:main_result}).  The associated functional equation
takes the irregular-graph reflection form $s\mapsto -s$ rather than
the regular-graph form $s\mapsto 1/(q s)$
(Section~\ref{sec:functional_equation}).  And the Ihara zeros are
algebraic, consistent with Paper~24's $\pi$-free structure
(Section~\ref{sec:closed_forms}).

We also address the natural companion question: could the Wick-rotated
side of the Fock projection --- scattering states on hyperbolic
$H^3$~\cite{bander1966} --- carry a Selberg zeta with proven RH analog?
The answer is no: GeoVac invariants supply no selection principle for
a discrete subgroup $\Gamma\subset\mathrm{SO}(3,1)$
(Section~\ref{sec:scope_boundary}).  This closes a door cleanly.  The
graph-zeta direction is structurally intrinsic; the Selberg direction
requires external arithmetic input that GeoVac does not provide.

% ===================================================================
\section{Ihara Zeta and Graph-RH}
\label{sec:ihara_background}
% ===================================================================

For a finite connected graph $G=(V,E)$ with $|V|$ vertices and $|E|$
edges, the \emph{Ihara zeta function}~\cite{ihara1966} is the formal
Euler product
\begin{equation}
\zeta_G(s) \;=\; \prod_{[C]}\,\bigl(1 - s^{\ell(C)}\bigr)^{-1},
\label{eq:ihara_euler}
\end{equation}
where $[C]$ ranges over equivalence classes of \emph{primitive} closed
walks in $G$ (walks with no backtracking and no tail, considered up to
cyclic rotation) and $\ell(C)$ is the walk length.  The product
converges for $|s|<1/\rho(T)$ where $T$ is the Hashimoto
non-backtracking edge operator (Section~\ref{sec:hashimoto}).

\subsection{Ihara--Bass determinantal formula}

A direct enumeration of primitive closed walks is infeasible for graphs
with more than a few dozen edges.  The Ihara--Bass
theorem~\cite{bass1992} replaces the Euler product by a finite-dimensional
determinant:
\begin{equation}
\zeta_G(s)^{-1} \;=\; \bigl(1 - s^2\bigr)^{r - c}\,
    \det\!\bigl(I - s\,A + s^2\,Q\bigr),
\label{eq:ihara_bass}
\end{equation}
where $A$ is the $|V|\times|V|$ adjacency matrix, $Q$ is the diagonal
matrix with entries $d_v - 1$ (where $d_v$ is the degree of vertex $v$),
$c$ is the number of connected components, and
$r = |E| - |V| + c = \beta_1$ is the first Betti number.  The
$(r-c)$ prefactor is the product-over-components generalization of
Terras's connected formula~\cite[Thm.~2.3]{terras2010}; it reduces to
the standard $(1-s^2)^{r-1}$ when $c=1$.

\subsection{Hashimoto non-backtracking operator}
\label{sec:hashimoto}

The \emph{Hashimoto edge operator}~\cite{hashimoto1989} $T$ acts on
the $2|E|$-dimensional space of directed edges of $G$: for each
directed edge $(u\to v)$,
\begin{equation}
T[(u\to v)] \;=\; \sum_{w\neq u,\; w\sim v} (v\to w).
\end{equation}
Equivalently $T$ is $|T|=2|E|\times 2|E|$, with entries
$T_{e,f}=1$ iff edge $f$ begins where edge $e$ ends and $f\neq e^{-1}$.
Then
\begin{equation}
\zeta_G(s)^{-1} \;=\; \det\!\bigl(I - s\,T\bigr),
\label{eq:ihara_hashimoto}
\end{equation}
so the reciprocal Ihara zeros of $G$ are exactly the nonzero
eigenvalues of $T$.

\subsection{Ramanujan graphs and graph-RH}

A $(q{+}1)$-regular graph is \emph{Ramanujan} if every non-trivial
eigenvalue $\lambda$ of the adjacency matrix satisfies
$|\lambda|\le 2\sqrt{q}$~\cite{lps1988,morgenstern1994}.  For irregular
graphs the Ramanujan condition is stated on the Hashimoto
spectrum: $G$ is \emph{(weakly) Ramanujan} if every non-trivial
eigenvalue $\mu$ of $T$ satisfies $|\mu|\le\sqrt{q_{\max}}$, where
$q_{\max} = \max_v(d_v - 1)$~\cite{kotani2000,terras2010}.
The non-trivial eigenvalues are those other than the Perron eigenvalue
$\rho(T)$ and its complex-conjugate reflection.

The \emph{graph Riemann Hypothesis} is the statement that $\zeta_G$ is
Ramanujan.  Equivalently, the Ihara zeros (reciprocals of $T$-eigenvalues)
are confined to the critical annulus
$1/\sqrt{q_{\max}} \le |s| \le 1$, outside of which only trivial zeros
(from the Perron eigenvalue) appear.  This is a \emph{proven}
theorem on specific graph families (the LPS-Margulis-Morgenstern
constructions~\cite{lps1988,morgenstern1994}) and has an extensive
literature; see Terras~\cite{terras2010} for a textbook account.

% ===================================================================
\section{The GeoVac Hopf Graphs}
\label{sec:graphs}
% ===================================================================

GeoVac builds discrete graphs on two natural geometries: the
three-sphere $S^3$ (the Fock~\cite{fock1935} projection for hydrogenic
bound states, Paper~7~\cite{paper7}) and the five-sphere $S^5$ (the
Bargmann--Segal projection for the 3D isotropic harmonic oscillator,
Paper~24~\cite{paper24}).  Both carry Hopf-fibration structure
($S^1\hookrightarrow S^3\to S^2$ and $S^1\hookrightarrow S^5\to
\mathbb{CP}^2$) and are the canonical discretizations referenced by
Paper~25~\cite{paper25}.  We summarize their combinatorial data.

\subsection{$S^3$ Coulomb graph}

Nodes are hydrogenic quantum-number triples $(n,\ell,m)$ with
$1\le n\le n_{\max}$, $0\le\ell<n$, $|m|\le\ell$.  Edges are generated
by the $\mathrm{SU}(2)\times\mathrm{SU}(1,1)$ ladder operators
$L_\pm, T_\pm$, which shift $m\mapsto m\pm 1$ (at fixed $n,\ell$) and
$n\mapsto n\pm 1$ (at fixed $\ell,m$), respectively.  The adjacency
rule therefore \emph{preserves}~$\ell$, which forces the graph to split
into disconnected components indexed by~$\ell$ (see
Section~\ref{sec:closed_forms}).

\subsection{$S^5$ Bargmann--Segal graph}

Nodes are SU(3) weights $(N,\ell,m_\ell)$ with $N\ge 0$,
$\ell\in\{N, N{-}2, N{-}4, \ldots, 0\text{ or }1\}$,
$|m_\ell|\le\ell$, truncated at $N\le N_{\max}$.  Edges are the $S^5$
dipole transitions with $\Delta N=\pm 1, \Delta\ell=\pm 1,
\Delta m_\ell\in\{0,\pm 1\}$, carrying exact-rational squared
matrix elements (Paper~24~\cite{paper24} \S III).  The adjacency is
\emph{not} $\ell$-preserving and the graph is connected at every
tested $N_{\max}$.

\subsection{Combinatorial summary}

\begin{table}[h]
\centering
\begin{tabular}{lrrrrr}
\toprule
Graph & $V$ & $E$ & $c$ & $\beta_1$ & $q_{\max}$ \\
\midrule
$K_4$ (sanity) & 4 & 6 & 1 & 3 & 2 \\
$K_{3,3}$ (sanity) & 6 & 9 & 1 & 4 & 2 \\
$S^3$ Coulomb, $n_{\max}=2$ & 5 & 3 & 2 & 0 & 1 \\
$S^3$ Coulomb, $n_{\max}=3$ & 14 & 13 & 3 & 2 & 2 \\
$S^5$ Bargmann--Segal, $N_{\max}=2$ & 10 & 15 & 1 & 6 & 4 \\
$S^5$ Bargmann--Segal, $N_{\max}=3$ & 20 & 42 & 1 & 23 & 8 \\
\bottomrule
\end{tabular}
\caption{Combinatorial data for the graphs tested.  $V$: vertices.
$E$: edges.  $c$: connected components.  $\beta_1=E-V+c$: first Betti
number.  $q_{\max}=\max_v(d_v-1)$: Ramanujan bound parameter.  The
$\beta_1=23$ at $S^5$ $N_{\max}=3$ agrees with the Paper~25 Hodge
computation of $\beta_1=110$ for the same graph family at $N_{\max}=5$
(the growth is roughly quadratic in $N_{\max}$).}
\label{tab:combinatorics}
\end{table}

% ===================================================================
\section{Main Result: the Ramanujan Property}
\label{sec:main_result}
% ===================================================================

Applying~(\ref{eq:ihara_bass}) and~(\ref{eq:ihara_hashimoto}) to each
graph gives the verdicts in Table~\ref{tab:ramanujan}.

\begin{table}[h]
\centering
\begin{tabular}{lrrrrr}
\toprule
Graph & $\rho(T)$ & $\max|\mu_{\mathrm{nontriv}}|$ & $\sqrt{q_{\max}}$ & Deviation & Ramanujan? \\
\midrule
$K_4$  & $2.0000$ & $1.4142$ & $1.4142$ & $\phantom{-}0.0000$ & YES (exact) \\
$K_{3,3}$ & $2.0000$ & $1.4142$ & $1.4142$ & $\phantom{-}0.0000$ & YES (exact) \\
$S^3$ C., $n_{\max}=2$ & $0$ & $0$ & $1.0000$ & $-1.0000$ & YES (trivial; forest) \\
$S^3$ C., $n_{\max}=3$ & $1.3532$ & $1.2157$ & $1.4142$ & $-0.1985$ & \textbf{YES} \\
$S^5$ BS, $N_{\max}=2$ & $2.3572$ & $1.7922$ & $2.0000$ & $-0.2078$ & \textbf{YES} \\
$S^5$ BS, $N_{\max}=3$ & $3.9053$ & $2.4716$ & $2.8284$ & $-0.3568$ & \textbf{YES} \\
\bottomrule
\end{tabular}
\caption{Ramanujan verdicts.  $\rho(T)$ is the Perron eigenvalue of the
Hashimoto operator (a trivial eigenvalue, not counted against the
Ramanujan bound).  $\max|\mu_{\mathrm{nontriv}}|$ is the largest
non-trivial Hashimoto eigenvalue.  Ramanujan requires
$\max|\mu_{\mathrm{nontriv}}|\le\sqrt{q_{\max}}$; the deviation
$\max|\mu_{\mathrm{nontriv}}|-\sqrt{q_{\max}}$ is strictly negative
for every non-trivial tested graph.}
\label{tab:ramanujan}
\end{table}

\begin{observation}
\label{obs:ramanujan}
For every GeoVac Hopf graph tested in this work, the largest
non-trivial Hashimoto eigenvalue is strictly smaller than the
Kotani--Sunada bound.  The graphs are therefore strictly
sub-Ramanujan: their spectral gap exceeds the optimal Ramanujan
gap of an equivalent-$q_{\max}$ regular graph.
\end{observation}

The strictness of the deviation is itself structurally interesting.
The Alon--Boppana bound~\cite{alon1986} asserts that for a sequence of
graphs of bounded degree, the largest non-trivial eigenvalue
asymptotes to the Ramanujan bound from below; graphs that achieve the
bound asymptotically are in this sense \emph{optimal}.  Our data at
three sizes ($S^3$, $S^5$ $N_{\max}=2$, $S^5$ $N_{\max}=3$) do not
suffice to establish the asymptotic behavior of the deviation -- we
explicitly leave this as an open question
(Section~\ref{sec:open_questions}).

% ===================================================================
\section{Closed-Form Ihara Zetas and Structural Factorizations}
\label{sec:closed_forms}
% ===================================================================

Working in exact rational arithmetic via \texttt{sympy}, we obtain
explicit closed forms for each graph's Ihara zeta.

\subsection{$S^3$ Coulomb at $n_{\max}=3$: per-$\ell$-shell
decomposition}

\begin{equation}
\zeta_G(s)^{-1} \;=\; -(s-1)^2 (s+1)^2 (s^2 - s + 1)(s^2 + s + 1)
    (2s^3 - s^2 + s - 1)(2s^3 + s^2 + s + 1).
\label{eq:s3_closed}
\end{equation}

The $S^3$ Coulomb graph at $n_{\max}=3$ has three connected components
indexed by $\ell=0,1,2$.  The $\ell=0$ component is a 3-node path
($\beta_1=0$, $\zeta_{\ell=0}=1$); the $\ell=2$ component is a 5-node
tree-like structure ($\beta_1=0$, $\zeta_{\ell=2}=1$).  \emph{All
non-trivial Ihara content lives in the $\ell=1$ component}, which is
the bipartite 3-prism / ladder graph $C_3\times P_2$.  The cyclotomic
factors $s^2 \pm s + 1$ in~(\ref{eq:s3_closed}) have zeros at primitive
sixth roots of unity (magnitude 1); they count the two triangles in
the 3-prism.  The irreducible cubics $2s^3 \pm s^2 + s \pm 1$ have
zeros at $|s|\approx 0.739$ (real) and $|s|\approx 0.823$ (complex),
corresponding to longer closed walks through the ladder.

The per-$\ell$ decomposition is a direct consequence of the Fock
adjacency preserving $\ell$.  It is an intrinsic structural property
of the $S^3$ Coulomb graph and one of the reasons the $S^3$ Ihara zeta
is analytically tractable.

\subsection{$S^5$ Bargmann--Segal at $N_{\max}=2$}

\begin{equation}
\zeta_G(s)^{-1} \;=\; -(s-1)^6 (s+1)^6 (2s^2 + 1)^2 (3s^4 + 3s^2 + 1)
    (24s^6 + 21s^4 + s^2 - 1).
\label{eq:s5_nmax2_closed}
\end{equation}

The graph is connected, $\beta_1=6$.  The $(s\pm 1)^6$ prefactor
captures the $\beta_1$ trivial zeros.  The remaining factors are all
even in $s$ and have integer coefficients.

\subsection{$S^5$ Bargmann--Segal at $N_{\max}=3$: the $12+22$
dichotomy}

\begin{equation}
\zeta_G(s)^{-1} \;=\; (s-1)^{23}(s+1)^{23}\cdot P_{12}(s) \cdot P_{22}(s),
\label{eq:s5_nmax3_closed}
\end{equation}
where
\begin{align}
P_{12}(s) &= 432 s^{12} + 666 s^{10} + 374 s^8 + 135 s^6 + 47 s^4 + 11 s^2 + 1, \\
P_{22}(s) &= 829440 s^{22} + 2453184 s^{20} + 3308104 s^{18} + 2696682 s^{16} \notag \\
           &\quad + 1470640 s^{14} + 557227 s^{12} + 146654 s^{10} + 25709 s^8 \notag \\
           &\quad + 2630 s^6 + 79 s^4 - 12 s^2 - 1.
\end{align}
Both $P_{12}$ and $P_{22}$ are polynomials in $s^2$ with integer
coefficients; they together account for all $12+22=34$ non-trivial
zeros, and combined with the $2\cdot 23 = 46$ trivial zeros at $s=\pm 1$
recover the total degree $2|E|=84$.

The $12+22$ split of the Ihara zeta is combinatorially nontrivial:
there is no obvious reason, from the adjacency alone, that the
characteristic polynomial of $T$ should factor into two irreducible
integer polynomials of these specific degrees.

\begin{observation}[$12+22$ factorization = $\mathbb{Z}_2$
$m_\ell$-reflection block decomposition]
\label{obs:12_22}
The two factors $P_{12}(s^2), P_{22}(s^2)$ in the $S^5$
Bargmann--Segal $N_{\max}=3$ Ihara zeta are exactly the characteristic
polynomials of the node-level Ihara--Bass matrix
$M(s) = I - sA + s^2 Q$ restricted to the two blocks of the
$\mathbb{Z}_2$ reflection $P:\; m_\ell\mapsto -m_\ell$.
Explicitly, $\det(M_{\text{sym}}) = (s-1)(s+1)\cdot P_{22}(s^2)$ on the
$13$-dimensional symmetric block (containing the $m_\ell=0$ fixed
points), and $\det(M_{\text{antisym}}) = P_{12}(s^2)$ on the
$7$-dimensional antisymmetric block.  The $\mathbb{Z}_2$ $m_\ell$-reflection
is the only sub-action of the Paper~25~\cite{paper25} Hopf $U(1)$ that
commutes with a real integer adjacency; the continuous $U(1)$ phase
rotation reduces to this $\mathbb{Z}_2$ when the data is real.
The analogous decomposition holds for the Dirac-$S^3$ Rule~B
$n_{\max}=3$ graph, with the $22+24$ dichotomy corresponding to the
equal-dimension $\mathbb{Z}_2$ blocks of the $m_j\to-m_j$
half-integer reflection (both blocks $14$-dimensional because $m_j$
is half-integer, so there are no reflection-fixed points).
\end{observation}

Observation~\ref{obs:12_22} was validated by direct restriction of
$M(s)$ to each reflection block in exact sympy rational arithmetic
(Sprint~2 Track RH-E, \texttt{debug/hopf\_u1\_block\_test\_memo.md}).
The product $\det(M_{\text{sym}})\cdot\det(M_{\text{antisym}})$ equals
the full $\det(M)$ symbolically.

\subsection{Algebraicity of zeros}

Every factor in every closed form above has integer coefficients.
The non-trivial Ihara zeros are therefore algebraic numbers, with
Galois conjugates also appearing in the zero set.  This is a direct
consequence of the adjacency matrix being $0/1$-valued and the
Ihara--Bass formula~(\ref{eq:ihara_bass}) being a determinant of an
integer-coefficient matrix polynomial.

\begin{corollary}[$\pi$-freeness of Ihara zeros]
\label{cor:pi_free}
No GeoVac Hopf Ihara zero has a closed form involving $\pi$, $\zeta(2)$,
$\zeta(3)$, or any other transcendental from the Paper~18 taxonomy.
Every zero is an algebraic number lying in a finite extension of
$\mathbb{Q}$.
\end{corollary}

Corollary~\ref{cor:pi_free} extends Paper~24's $\pi$-free
certificate from adjacency data to Ihara-zeta data.  The same
bit-exact $\pi$-freeness is preserved under Ihara-zeta.

% ===================================================================
\section{Functional Equation: the Irregular-Graph Reflection}
\label{sec:functional_equation}
% ===================================================================

For a $(q{+}1)$-regular graph, the Ihara zeta satisfies the
Stark--Terras functional equation~\cite{stark1996,terras2010}
\begin{equation}
\Xi_G(s) := (1 - s^2)^{r-1}\bigl(1 - (qs)^2\bigr)^{r-1}\zeta_G(s)^{-1},
\quad \Xi_G(s) = \Xi_G(1/(qs)).
\label{eq:stark_terras}
\end{equation}
In particular the Ihara zeros lie on the critical circle $|s|=1/\sqrt{q}$
when the graph is Ramanujan.  Both sanity witnesses ($K_4, K_{3,3}$)
satisfy~(\ref{eq:stark_terras}) with zeros on $|s|\in\{\tfrac{1}{\sqrt{2}},
1\}$.

\emph{Neither} GeoVac Hopf graph is regular: $q$ varies from 0 (at
leaf nodes) to $q_{\max}$ (at dense interior nodes).  The Stark--Terras
single-circle statement therefore does not hold; instead, the
non-trivial Ihara zeros populate an \emph{annulus}
$1/\sqrt{q_{\max}}\le|s|\le 1$, and their distribution within the
annulus is constrained only by the Hashimoto spectrum.

What \emph{does} hold uniformly across our graphs:

\begin{proposition}[Reflection symmetry]
\label{prop:reflection}
For each GeoVac Hopf graph tested, $\zeta_G(s)^{-1}$ is invariant
under $s\mapsto -s$ up to a possible sign factor of the trivial
$(s\pm 1)^{\beta_1}$ prefactors.  Equivalently, every non-trivial
polynomial factor in the closed forms
(\ref{eq:s3_closed}), (\ref{eq:s5_nmax2_closed}),
(\ref{eq:s5_nmax3_closed}) is even in~$s$.
\end{proposition}

The proposition is a structural statement about undirected graphs:
the Hashimoto matrix $T$ satisfies $T\sim -T$ under the involution
$e\mapsto e^{-1}$ on directed edges, so its characteristic polynomial
is even in its argument.  Proposition~\ref{prop:reflection} is
therefore a universal property of the Ihara zeta on any undirected
graph, regular or not.

\begin{proposition}[Algebraic closure under Galois conjugation]
\label{prop:galois}
The non-trivial Ihara zeros of every GeoVac Hopf graph are algebraic
numbers whose Galois conjugates also appear in the zero set.  The
zero set is therefore a finite union of Galois orbits in
$\bar{\mathbb{Q}}$.
\end{proposition}

This is immediate from the integer-coefficient factorizations of
Section~\ref{sec:closed_forms}.

The weakening from the full Stark--Terras symmetry to
Propositions~\ref{prop:reflection} and~\ref{prop:galois} is a feature
of irregularity, not a pathology of our graphs.  It is
Kotani--Sunada's~\cite{kotani2000} Theorem~1.3 which supplies the
correct $\Xi$-completion for the irregular case: a product over all
degree strata of the graph, rather than a single regular-graph
correction.  Making that completion explicit for our two graphs is a
useful computational exercise but is not required to state the
Ramanujan property, which is already secured by
Observation~\ref{obs:ramanujan}.

% ===================================================================
\section{Scope Boundary: No Natural Hyperbolic Quotient}
\label{sec:scope_boundary}
% ===================================================================

A natural companion question to the Ihara-zeta result is whether the
Wick-rotated continuation of the Fock projection, which takes the
compact $S^3$ to the non-compact hyperbolic 3-space $H^3$ at positive
energy~\cite{bander1966}, delivers a Selberg zeta with a \emph{proven}
RH analog (on congruence Bianchi quotients~\cite{egm1998}).  We
investigated this in parallel and report a clean negative.

\subsection{The Wick rotation is clean but continuous}

Following Bander--Itzykson~\cite{bander1966}, the Fock projection with
$p_0=\sqrt{-2mE}$ continues to $p_0=i\kappa_{\mathrm{phys}}$ at
$E>0$, rotating the target manifold from the round $S^3\subset
\mathbb{R}^4$ to the upper unit pseudosphere in $\mathbb{R}^{3,1}$,
i.e., hyperbolic 3-space $H^3$ in its Minkowski embedding.  The
symmetry group rotates from compact $\mathrm{SO}(4)$ to non-compact
$\mathrm{SO}(3,1)$, with the SO(4) principal-series
representations continuing to the principal/complementary series of
$\mathrm{SO}(3,1)$.  The Coulomb scattering amplitudes are matrix
elements in this continuous principal series.  This is clean and
unambiguous.

\subsection{But no discrete quotient $\Gamma\subset\mathrm{SO}(3,1)$
is singled out}

To recover a graph-like discrete object on $H^3$, one needs a
discrete subgroup $\Gamma$ with quotient $H^3/\Gamma$.  Of the four
candidate $\Gamma$'s that GeoVac structure suggests, none works:

\begin{enumerate}
\item \textbf{Weyl group of $\mathrm{SO}(4)$.}  The Weyl group
$W(D_2)=\mathbb{Z}/2\times\mathbb{Z}/2$ has order 4.  Finite, not a
lattice in $\mathrm{SO}(3,1)$.  Acting on $H^3$ it has fixed points
everywhere; it does not produce a discrete hyperbolic quotient.
\item \textbf{Hopf $S^1$ fiber.}  The compact Hopf $S^1$ is
structurally irrelevant for $H^3$: hyperbolic 3-space has no Hopf-type
principal bundle with a compact $S^1$ fiber.  This is the same
Coulomb/HO asymmetry Paper~24~\cite{paper24} identified, one layer
further.
\item \textbf{Bianchi groups $\mathrm{PSL}(2,\mathcal{O}_K)$,
$K=\mathbb{Q}(\sqrt{-d})$.}  The class-number-one cases
$d\in\{1,2,3,7,11\}$ act discontinuously on $H^3$.  None of GeoVac's
numerical invariants -- the Rydberg constant $R_\infty$, the Fock
momentum $p_0=Z/n$, $\kappa=-\tfrac{1}{16}$, $B=42$, $F=\pi^2/6$,
$\Delta=\tfrac{1}{40}$, $\alpha^{-1}\approx 137$ -- matches any
specific $d$.  The Bianchi groups are objects of number theory, not
of hydrogen.
\item \textbf{Picard group $\mathrm{PSL}(2,\mathbb{Z}[i])$.}  As the
simplest Bianchi case, Picard is the default.  But ``simplest'' is
not a physics selection principle, and the Coulomb node labels
$(n,\ell,m)\in\mathbb{Z}^3$ are not in $\mathbb{Z}[i]$.  No natural
selection.
\end{enumerate}

The literature survey~\cite{rh_survey} that accompanied this
investigation confirmed that \emph{no published work} links the
$\mathrm{SO}(4)$ Coulomb symmetry to Bianchi arithmetic quotients.

\begin{observation}[No natural hydrogen-Selberg bridge]
\label{obs:no_selberg}
The Fock Wick rotation $S^3\to H^3$ is well-defined and continuous,
but GeoVac provides no selection principle for a discrete subgroup
$\Gamma\subset\mathrm{SO}(3,1)$.  The Selberg-zeta angle on
hydrogen-derived hyperbolic quotients is therefore structurally
closed.
\end{observation}

This observation refines the Paper~24~\cite{paper24} universal /
Coulomb-specific partition one level further:
\begin{itemize}
\item $S^3$ (Coulomb, bound): compact, Hopf, calibration $\pi$, discrete
spectrum, graph-zeta available.
\item $S^5$ (HO, Bargmann--Segal): compact complex, Hopf $U(1)$,
bit-exactly $\pi$-free, discrete spectrum, graph-zeta available.
\item $H^3$ (Coulomb, continuum): non-compact, no natural Hopf
bundle, no natural discrete subgroup, continuous spectrum,
Selberg-zeta \emph{requires external arithmetic input}.
\end{itemize}

The graph-zeta direction is therefore the single RH-adjacent direction
that is intrinsic to GeoVac.  The Selberg direction, for hydrogen,
is not.

% ===================================================================
\section{Open Questions}
\label{sec:open_questions}
% ===================================================================

\paragraph{Extension to $N_{\max}=5$.}  The Paper~25 headline case
for the $S^5$ Bargmann--Segal graph reports $V=56$, $E=165$,
$\beta_1=110$.  The Hashimoto operator at that size is $2|E|=330$-dimensional
-- a fully tractable numerical eigensolve.  Extending the Ramanujan
verdict and the closed-form factorization to $N_{\max}=5$ would give
three data points on the deviation's growth (already at $-0.357$
for $N_{\max}=3$) and a headline match to Paper~25's Hodge data.

\paragraph{Alon--Boppana asymptotics --- resolved (Sprint~2, RH-D).}
Extending the Ramanujan verdict to $n_{\max},N_{\max}\in\{4,5\}$ in a
post-draft sweep establishes that the signed deviation grows
\emph{linearly in $V$} for three of four families (scalar $S^3$
Coulomb: $R^2=0.994$; scalar $S^5$ Bargmann--Segal: $R^2=0.766$;
Dirac-$S^3$ Rule~B: $R^2=0.828$).  The Kotani--Sunada bound is crossed
between $V\approx 30$ and $V\approx 60$ for each of these families:
scalar $S^3$ at $n_{\max}=5$ ($V=55$, deviation $+0.20$), scalar $S^5$
at $N_{\max}=5$ ($V=56$, deviation $+0.65$), Dirac-$S^3$ Rule~B at
$n_{\max}=4$ ($V=60$, deviation $+1.53$).  Only the Dirac-$S^3$ Rule~A
family appears to remain sub-Ramanujan, with a $1/V$-class envelope
($|\text{deviation}|\sim V^{-3.05}$, $R^2=0.994$).  The conjectured
structural mechanism is \emph{mixed-density topology}: $q_{\max}$ is
set by the single densest vertex while the Perron eigenvalue of $T$ is
driven by graph volume; as the graph grows, near-regular dense
sub-blocks appear whose Perron exceeds $\sqrt{q_{\max}}$.  The
graph-RH for GeoVac is a finite-size statement, secured for
$n_{\max},N_{\max}\le 3$ and failing asymptotically for the three
dense families.  Full data at \texttt{debug/data/alon\_boppana\_sweep.json}.

\paragraph{Hopf-quotient block decomposition --- validated (Sprint~2, RH-E).}
Observation~\ref{obs:12_22} has been verified by direct restriction of
the node-level Ihara--Bass matrix $M(s)$ to each $\mathbb{Z}_2$
$m$-reflection block in exact sympy rational arithmetic.  The block
decomposition reproduces $P_{12}$ and $P_{22}$ (scalar $S^5$
$N_{\max}=3$) and $P_{22}$ and $P_{24}$ (Dirac-$S^3$ Rule~B
$n_{\max}=3$) symbolically.  The $\mathbb{Z}_2$ reflection is the only
sub-action of the Paper~25 Hopf $U(1)$ that commutes with a real
integer adjacency, and the mechanism is identical in scalar and spinor
cases.

\paragraph{Spin-ful Ihara zeta on the Dirac-$S^3$ graph.}  The
infrastructure built for the Tier-3 relativistic
sprint~\cite{paper14,paper18} -- the Dirac spectrum $|\lambda_n|=n+\tfrac32$,
the $(\kappa, m_j)$-labeled spinor harmonic basis, the
Szmytkowski matrix elements -- induces a distinct graph structure
on the Dirac sector.  Its Ihara zeta is a genuinely new object:
the edges differ from the Schr\"odinger graph (the Dirac rule
couples spin and orbital angular momentum), and the resulting Hashimoto
spectrum is -- conjecturally -- richer.  This is the natural spinor
extension of the present paper and the recommended next Ihara-zeta
sprint from both the Wick memo and the present work.

\paragraph{Continuum limit --- redirected from Ihara to spectral zeta
(Sprint~2, RH-G).}  Sprint~2 Track~RH-G investigated whether any
combinatorial or refinement limit of the GeoVac Ihara zeta connects
to a classical $L$-function.  The headline obstruction is structural:
the GeoVac Hopf graphs are \emph{bipartite} (all closed
non-backtracking walks have even length), so $\operatorname{tr}(T^L)=0$
for odd $L$ and the Ihara walk-length spectrum is supported on the
even integers.  Paper~28's $\chi_{-4}$ Dirichlet character has support
only on odd integers, so the $\chi_{-4}$-twisted Ihara zeta
$\zeta_G(s;\chi_{-4}) = \prod_{[C]}(1 - \chi_{-4}(\ell(C))s^{\ell(C)})^{-1}$
is identically the trivial Euler product.  The $\chi_{-4}$ structure
that does appear in Paper 28 lives on the \emph{spectral} side --- the
Dirac Dirichlet series $D(s) = \sum_{n} g_n\cdot |\lambda_n|^{-s}$ at
quarter-integer Hurwitz shifts, splitting into $D_{\text{even}}(s)$ and
$D_{\text{odd}}(s)$ whose differences expose Catalan's $G=\beta(2)$
and $\beta(4)$.  Supporting statistics: the Hashimoto eigenvalue-spacing
coefficient of variation is $\approx 1$ (Poisson, Berry--Tabor), not
the GUE value $\approx 0.42$, categorically excluding a direct
Hilbert--P\'olya identification with $\zeta$ zeros.  The bridge to
classical $\zeta(s)$, if any, must therefore go through the
\emph{spectral} zeta (Paper~28's machinery), not through the Ihara
zeta.  Full obstruction analysis at
\texttt{debug/riemann\_limit\_memo.md}.

% ===================================================================
\section{Conclusion}
\label{sec:conclusion}
% ===================================================================

The GeoVac Hopf graphs satisfy the graph Riemann Hypothesis --- the
Kotani--Sunada Ramanujan bound --- strictly, at every small size
tested ($n_{\max},N_{\max}\le 3$).  Extending to $n_{\max},N_{\max}
\in\{4,5\}$ in a Sprint-2 follow-up sweep shows that three of four
graph families cross the bound at modest sizes
($V\approx 30$--$60$) with a signed deviation growing linearly in $V$,
so the graph-RH is a finite-size statement for GeoVac, not an
asymptotic one (Section~\ref{sec:open_questions}, Alon--Boppana
subsection).  Their Ihara zetas factor with integer coefficients; their
zeros are algebraic; no transcendentals enter the zero set, consistent
with Paper~24's $\pi$-free certificate.  The per-$\ell$-shell
decomposition of the $S^3$ Coulomb zeta and the $12+22$ / $22+24$
dichotomies of the scalar $S^5$ $N_{\max}=3$ and Dirac-$S^3$ Rule~B
$n_{\max}=3$ zetas are, by direct verification, the $\mathbb{Z}_2$
$m$-reflection blocks of the node-level Ihara--Bass matrix
(Observation~\ref{obs:12_22}); these connect back to Paper~7's SO(4)
selection rules and to the real-integer-adjacency-compatible
sub-action of Paper~25's Hopf-$U(1)$ gauge structure.  Neither graph
is regular, so the single-critical-circle Stark--Terras functional
equation does not apply; what does hold uniformly is the
$s\mapsto -s$ reflection and Galois closure of the zero set.

The adjacent Selberg-zeta direction -- a Wick rotation of the Fock
projection to hyperbolic 3-space, followed by a quotient by a
discrete arithmetic subgroup -- is structurally closed for hydrogen:
the framework supplies no selection principle for such a subgroup.
The graph-zeta angle is therefore the single RH-adjacent direction
intrinsic to GeoVac.

We frame this as an observation rather than a claim about the
classical Riemann Hypothesis.  The graph-RH for a specific family
of finite graphs is a different object than RH for $\zeta(s)$.  What
these results establish is a concrete bridge at small graph sizes:
\emph{our} framework has an Euler product, a functional equation (in
the irregular-graph form), a critical region, and zeros that respect
the Ramanujan bound at $n_{\max},N_{\max}\le 3$, plus a validated
$\mathbb{Z}_2$ block-decomposition matching the Paper~25 Hopf
structure.  The Sprint~2 search for a limiting bridge to classical
$\zeta(s)$ (Track~RH-G) produced a structural redirection: the
bipartiteness of the GeoVac Hopf graphs blocks any Ihara-side bridge,
but Paper~28's spectral Dirichlet series (equipped with
$\chi_{-4}$) remains the open direction on the spectral side.  The
Ihara zeta therefore terminates as a self-contained structural object
with validated $\mathbb{Z}_2$ block decomposition, and the Prize-bound
research program moves to the spectral-zeta side.

The computations reported here are recorded in
\texttt{geovac/ihara\_zeta.py} and \texttt{tests/test\_ihara\_zeta.py},
with full numerical data in
\texttt{debug/data/ihara\_zeta\_geovac\_hopf.json} and the
supporting Wick-rotation memo at
\texttt{debug/fock\_continuation\_memo.md}.  All Ihara zeta closed
forms above are exact in \texttt{sympy} rational arithmetic.

% ===================================================================
\begin{thebibliography}{99}
% ===================================================================

\bibitem{paper1} GeoVac Project.  \emph{Paper 1: Spectral Graph Theory
and the Geometric Vacuum.}  GeoVac Project, ongoing.
\bibitem{paper7} GeoVac Project.  \emph{Paper 7: The Dimensionless
Vacuum.}  GeoVac Project, ongoing.  18 symbolic proofs of the
$S^3$ conformal equivalence.
\bibitem{paper14} GeoVac Project.  \emph{Paper 14: Qubit Encoding of
the Geometric Vacuum.}  GeoVac Project, ongoing.
\bibitem{paper18} GeoVac Project.  \emph{Paper 18: Exchange Constants
and Transcendental Taxonomy.}  GeoVac Project, ongoing.
\bibitem{paper24} GeoVac Project.  \emph{Paper 24: Bargmann--Segal
Lattice: a $\pi$-free Discretization of the 3D Harmonic Oscillator.}
GeoVac Project, 2026.
\bibitem{paper25} GeoVac Project.  \emph{Paper 25: The GeoVac Hopf
Graph as a Lattice-Gauge Structure.}  GeoVac Project, April 2026.
\bibitem{paper28} GeoVac Project.  \emph{Paper 28: QED on $S^3$:
Spectral Geometry of Perturbative Transcendentals.}  GeoVac Project,
April 2026.
\bibitem{fock1935} V. Fock.  Zur Theorie des Wasserstoffatoms.
\emph{Z.\ Phys.}\ \textbf{98}, 145 (1935).
\bibitem{bander1966} M. Bander and C. Itzykson.  Group Theory of the
Hydrogen Atom.  \emph{Rev.\ Mod.\ Phys.}\ \textbf{38}, 330 (1966).
\bibitem{ihara1966} Y. Ihara.  Discrete subgroups of $\mathrm{PL}(2,
\mathbb{Q}_p)$.  \emph{J.\ Math.\ Soc.\ Japan}\ \textbf{18}, 219 (1966).
\bibitem{bass1992} H. Bass.  The Ihara-Selberg zeta function of a
tree lattice.  \emph{Internat.\ J.\ Math.}\ \textbf{3}, 717 (1992).
\bibitem{hashimoto1989} K.-i. Hashimoto.  Zeta functions of finite
graphs and representations of $p$-adic groups.  Advanced Studies in
Pure Math.\ \textbf{15}, 211 (1989).
\bibitem{kotani2000} M. Kotani and T. Sunada.  Zeta functions of
finite graphs.  \emph{J.\ Math.\ Sci.\ Univ.\ Tokyo}\ \textbf{7}, 7
(2000).
\bibitem{terras2010} A. Terras.  \emph{Zeta Functions of Graphs:\
A Stroll Through the Garden.}  Cambridge Univ.\ Press, 2010.
\bibitem{lps1988} A. Lubotzky, R. Phillips, and P. Sarnak.  Ramanujan
graphs.  \emph{Combinatorica}\ \textbf{8}, 261 (1988).
\bibitem{morgenstern1994} M. Morgenstern.  Existence and explicit
constructions of $(q+1)$-regular Ramanujan graphs for every prime
power $q$.  \emph{J.\ Combin.\ Theory Ser.\ B}\ \textbf{62}, 44 (1994).
\bibitem{stark1996} H. Stark and A. Terras.  Zeta functions of finite
graphs and coverings.  \emph{Adv.\ Math.}\ \textbf{121}, 124 (1996).
\bibitem{alon1986} N. Alon and V. Milman.  $\lambda_1$, isoperimetric
inequalities for graphs and superconcentrators.  \emph{J.\ Combin.\
Theory Ser.\ B}\ \textbf{38}, 73 (1985); N. Alon.  Eigenvalues and
expanders.  \emph{Combinatorica}\ \textbf{6}, 83 (1986).
\bibitem{egm1998} J. Elstrodt, F. Grunewald, and J. Mennicke.
\emph{Groups Acting on Hyperbolic Space: Harmonic Analysis and Number
Theory.}  Springer, 1998.
\bibitem{rh_survey} GeoVac Project.  Riemann-Hypothesis literature
survey, \texttt{debug/rh\_literature\_survey.md}, April 2026.

\bibitem{matsuura2025} S. Matsuura and K. Ohta.  Kazakov-Migdal model
on graphs and the Ihara zeta function.  \emph{Prog.\ Theor.\ Exp.\
Phys.}\ \textbf{2025}, 063B01 (2025);
arXiv:2204.06424 (2022).  Realises the Ihara zeta as a partition
function of a lattice gauge theory on a generic finite graph; the
closest published precedent to treating our Hopf-structured graphs
as carriers of a graph zeta, although it does not cover spherical /
Fock-projected graphs specifically.

\bibitem{yakaboylu2024} E. Yakaboylu.  Hermitian reflection of the
Bender--Brody--M\"uller operator and the Riemann zeros.
\emph{J.\ Phys.\ A: Math.\ Theor.}\ \textbf{57}, 235204 (2024);
arXiv:2408.15135.  A current-generation Hilbert--P\'olya candidate;
the operator lives on $L^2([0,\infty))$, a half-line and not a
compact manifold, so it does not directly apply to a GeoVac graph
and reinforces the observation in Section~\ref{sec:scope_boundary}
that the compact-manifold analogues of this program remain open.

\bibitem{hmy2024} J. Huang, B. McKenzie, and H.-T. Yau.  Ramanujan
property and edge universality of random regular graphs.
arXiv:2412.20263 (2024).  Recent progress on the probabilistic
Ramanujan property for random regular graphs; useful as ambient
context for sub-Ramanujan deviation behaviour observed in
Section~\ref{sec:main_result}.

\end{thebibliography}
\end{document}
